% =============================================================================
% Formatação central do artigo. Fonte das regras: esquema transcrito pelos autores
% em 09/11/2026; as divergências com o guia citado e as verificações possíveis
% estão em article/compliance_report.md. O modelo do AVA não foi localizado.
% =============================================================================

% ---------------------------------------------------------------------------
% Página e margens (A4): superior e esquerda 3 cm; inferior e direita 2 cm
% ---------------------------------------------------------------------------
\geometry{
  a4paper,
  left=30mm,
  right=20mm,
  top=30mm,
  bottom=20mm,
  headheight=12pt,
  headsep=6mm,      % número da página a ~2 cm da borda superior
  footskip=10mm,
}

% ---------------------------------------------------------------------------
% Fonte e espaçamento: Arial 12; CORPO com espaçamento 1,5 (instrução expressa dos autores).
% 1,5 linha do Word com Arial 12 = 1,5 x 13,8 pt = 20,7 pt = 1,428 x 14,5 pt (linha-base do LaTeX a 12 pt).
% Única definição do espaçamento do corpo: \espacamentocorpo. Não use \setstretch/\linespread locais.
% Exceções (espaçamento simples, definidas só aqui): tabelas, notas de fonte, pseudocódigo,
% legendas, referências e qualificações. Parágrafos sem recuo, com 12 pt antes.
% ---------------------------------------------------------------------------
\renewcommand{\familydefault}{\sfdefault}
\newcommand{\espacamentocorpo}{1.428}
\setstretch{\espacamentocorpo}
\setlength{\parindent}{0pt}
\setlength{\parskip}{12pt}

\raggedbottom
\setlength{\emergencystretch}{1.5em}  % evita linhas que excedem a margem perto de citações longas
\clubpenalty=10000
\widowpenalty=10000
\displaywidowpenalty=10000

% Tamanhos auxiliares
\newcommand{\fontedez}{\setstretch{1}\fontsize{10}{11.5}\selectfont}
\newcommand{\fonteonze}{\setstretch{1}\fontsize{11}{12.65}\selectfont}
\newcommand{\fontedozesimples}{\setstretch{1}\fontsize{12}{13.8}\selectfont}
% Conteúdo tabular compacto, com folga vertical junto às linhas da grade.
\newcommand{\fontetabela}{\setstretch{1}\fontsize{9.5}{11.5}\selectfont\setlength{\tabcolsep}{3pt}\setlength{\extrarowheight}{1pt}\renewcommand{\arraystretch}{1.18}}

% ---------------------------------------------------------------------------
% Títulos de seção sem numeração, alinhados à esquerda
% ---------------------------------------------------------------------------
\setcounter{secnumdepth}{0}
\titleformat{\section}{\normalfont\raggedright\fontsize{12}{20.7}\selectfont\bfseries}{\thesection}{0.5em}{\MakeUppercase}
\titleformat{\subsection}{\normalfont\raggedright\fontsize{12}{20.7}\selectfont}{\thesubsection}{0.5em}{\MakeUppercase}
\titleformat{\subsubsection}{\normalfont\raggedright\fontsize{12}{20.7}\selectfont}{\thesubsubsection}{0.5em}{\MakeUppercase}
\titlespacing*{\section}{0pt}{12pt}{0pt}
\titlespacing*{\subsection}{0pt}{0pt}{0pt}
\titlespacing*{\subsubsection}{0pt}{0pt}{0pt}

% ---------------------------------------------------------------------------
% Número da página: algarismo arábico, Arial 10, canto superior direito
% ---------------------------------------------------------------------------
\fancypagestyle{artigo}{%
  \fancyhf{}%
  \fancyhead[R]{\fontsize{10}{12}\selectfont\thepage}%
  \renewcommand{\headrulewidth}{0pt}%
  \renewcommand{\footrulewidth}{0pt}%
}
\pagestyle{artigo}

% ---------------------------------------------------------------------------
% Listas (alíneas)
% ---------------------------------------------------------------------------
\setlist{
  labelindent=6.4mm,
  leftmargin=12.7mm,
  labelwidth=6.3mm,
  labelsep=0pt,
  align=left,
  topsep=0pt,
  partopsep=0pt,
  itemsep=4pt,
  parsep=0pt,
}
\setlist[itemize,1]{label=\textbullet}
\setlist[enumerate,1]{label=\arabic*.}

% ---------------------------------------------------------------------------
% Figuras, tabelas e quadros: título acima, Arial 11; fonte abaixo, Arial 11.
% ---------------------------------------------------------------------------
\DeclareCaptionLabelSeparator{hifen}{ - }
\DeclareCaptionFont{onze}{\fontsize{11}{12.65}\selectfont\setstretch{1}}  % legendas em espaçamento simples
\captionsetup{
  font=onze,
  labelfont=normalfont,
  textfont=normalfont,
  labelsep=hifen,
  justification=centering,
  singlelinecheck=true,
  position=top,
  skip=6pt,
}

\floatplacement{figure}{htbp}
\floatplacement{table}{htbp}

\newcommand{\largurapadraofigura}{0.6\linewidth}
\newcommand{\alturamaximafigura}{0.4\textheight}
\newcommand{\espacofonte}{6pt}

% Notas de rodapé: Arial 10, entrelinha simples, alinhadas à esquerda.
\let\footnotesize\fontedez
\renewcommand{\footnoterule}{\kern-3pt\hrule width 40mm\kern 2.6pt}

\setlength{\intextsep}{14pt plus 2pt minus 2pt}
\setlength{\textfloatsep}{10pt plus 2pt minus 2pt}
\setlength{\floatsep}{10pt plus 2pt minus 2pt}

% Citação direta longa: Arial 10, simples, recuo de 4 cm
\newcommand{\recuocitacao}{40mm}

% Referências: Arial 12 simples, alinhadas à esquerda, sem recuo
\newcommand{\espacoentrereferencias}{6pt}

% Parâmetros de flutuantes: evitam páginas de figuras com espaço vertical esticado
\makeatletter
\setlength{\@fptop}{0pt}
\setlength{\@fpsep}{12pt}
\setlength{\@fpbot}{0pt plus 1fil}
\makeatother
\renewcommand{\topfraction}{0.85}
\renewcommand{\bottomfraction}{0.7}
\renewcommand{\textfraction}{0.1}
\renewcommand{\floatpagefraction}{0.9}
\setcounter{topnumber}{3}
\setcounter{totalnumber}{4}
